\section{The recursive symmetric-group family}
\label{sec:family}
The smaller family in \cite[Example 5.2]{DMV} contains explicit information
not implied by the bare existence theorem. We retain its construction, color
counts and faithful top action.

\begin{theorem}\label{R.SymmetricFamily}
For each $n\geq2$ there is a connected properly colored triangle-free graph
$X_n$, with colors $0,\ldots,n$, satisfying
\[
 (C(X_n),K(X_n))\cong(S_n,A_n).
\]
Its color sizes are
\begin{equation}\label{eq:class-sizes}
 |t^{-1}(n)|=n,\qquad
 |t^{-1}(i)|=\frac{n!}{2(i-1)!}\quad(0\leq i<n),
\end{equation}
where $(-1)!=0!=1$. The top class carries the faithful natural $S_n$-action.
An odd permutation exchanges exactly colors zero and one, while an even
permutation fixes every color. The minimum degree is at least two for $n\geq3$.
\end{theorem}
\begin{proof}
Let $X_2$ be the four-vertex path with successive colors $2,0,1,2$.
Label its endpoints $1,2$. Its only correlations are the identity and the
reversal; reversal exchanges colors zero and one. This gives the claimed
pair and faithful top action. The endpoints have degree one, so this is a
pre-induction object, not a minimum-degree-two base case.

Suppose $X_{n-1}$ has been constructed, with its top vertices labeled
$1,\ldots,n-1$. Introduce vertices $v_1,\ldots,v_n$ of color $n$ and
vertices $v_{\{i,j\}}$ of color $n-1$, joined to $v_i$ and $v_j$.
Choose even permutations $g_j\in A_n$ with $g_j(n)=j$ and $g_n=1$;
these exist because $A_n$ is transitive for $n\geq3$.
At each $j$ attach a copy of $X_{n-1}$, identifying its top vertex labeled
$i$ with $v_{\{g_j(i),j\}}$. The interiors of different copies are disjoint,
and all their lower colors are inherited.

The graph is connected: each copy is connected and meets its boundary, and
the new top and pair vertices form the subdivided complete graph.
Within every copy the coloring is proper and there is no triangle.
Different copies meet only at pair vertices, all of the same color and hence
pairwise nonadjacent. The neighbors of a new top vertex also all have color
$n-1$ and are pairwise nonadjacent. Thus no triangle is introduced.
An interior vertex retains its old degree. A top vertex has degree $n-1$;
a pair vertex gets its two new top neighbors and its neighbors in two copies.
For $n=3$ only pair vertices had been endpoints of the base paths: their
new degree is four, whereas the other vertices have degree two.
For larger $n$ the minimum-degree assertion follows by induction.

There are $n$ top vertices and $\binom n2$ pair vertices. For $i<n-1$,
the class of color $i$ is the disjoint union of $n$ copies of the previous
class. These recurrences and the base path prove~\eqref{eq:class-sizes}.

For $\sigma\in S_n$ set $v_i\mapsto v_{\sigma(i)}$ and
$v_{\{i,j\}}\mapsto v_{\{\sigma(i),\sigma(j)\}}$. On the copy at $j$
use the inductively constructed correlation corresponding to
\[
 h_{\sigma,j}=g_{\sigma(j)}^{-1}\sigma g_j.
\]
This permutation fixes $n$ and therefore acts on the $n-1$ labels of that
copy. Its sign equals the sign of $\sigma$, because the chosen representatives
are even. The induction hypothesis says that all these local maps have the
same palette permutation: the identity for even $\sigma$, or exchange of
zero and one for odd $\sigma$.

On a boundary vertex labeled $i$ in the copy at $j$, the image has global
label
\[
 \{g_{\sigma(j)}h_{\sigma,j}(i),\sigma(j)\}
   =\{\sigma g_j(i),\sigma(j)\}.
\]
Thus local maps agree with the displayed map on every shared pair vertex.
They combine into a correlation, with inverse given by $\sigma^{-1}$.
Moreover
\[
 h_{\tau,\sigma(j)}h_{\sigma,j}=h_{\tau\sigma,j},
\]
so their composition respects multiplication. We have constructed an
$S_n$-action with the stated palette kernel and natural action on the top.

It remains to exclude extra correlations. For $n\geq4$, the top class of
size $n$ is uniquely smallest, since the next size is $\binom n2>n$ and all
lower classes are larger. For $n=3$ all four classes have size three, but
the top vertices are precisely the degree-two vertices both of whose
neighbors have degree four. Each other degree-two vertex has one neighbor
of degree two and one of degree four. Therefore every correlation preserves
the top class, in every case.

A correlation fixing the top pointwise fixes each pair vertex, since its
two top neighbors characterize it. We verify the component-recognition
hypothesis needed before applying induction to a copy. The graph $X_m$
with its top removed is connected for every $m\geq2$. For $m=2$ it is an
edge. For $m\geq3$, it is the union of its $m$ connected attached copies,
and any two copies meet in their common pair vertex. Hence it is connected.

In $X_n$, deleting both top and pair vertices consequently leaves exactly
the connected interiors of the attached copies. The boundary of the interior
at $j$ is
\[
 \{v_{\{i,j\}}:i\ne j\}.
\]
Every boundary vertex meets the interior: in the base it is an endpoint of
a path, and at every subsequent stage each top vertex has neighbors in the
lower part. These boundary stars are distinct for $n\geq3$. A correlation
fixing all pair vertices must therefore preserve each copy. On that copy
it is a correlation fixing the top pointwise, and is the identity by
inductive faithfulness. This is the component-boundary rigidity argument
of the common paper.

Thus restriction embeds the full correlation group in $S_n$. The constructed
action already realizes every member of $S_n$, and its palette kernel is
$A_n$. This proves all claims.
\end{proof}

For $n\geq3$, the common chamber theorem immediately produces
$(S_n,A_n)$-geometries of rank $n+1$. For $n=2$, first clean the base path
and then complete it, or use Theorem~\ref{R.FinitePairs}. The source's
blanket minimum-degree claim must be qualified at this base: it is not
legitimate to apply the chamber theorem directly to its degree-one endpoints.
